\chapter{Sequences}

\section{Motivation and definition}

I'm sure that you are probably already familiar with the word \textbf{sequence}. $1, 2, 3, 4, 5, \dots$ is a sequence of numbers which are increasing by one, $2, 4, 8, 16, 32, \dots$ is a sequence of numbers which are doubling. If the sequence is increasing by $d$ each time, then it's an arithmetic sequence. If it's multiplied by $r$ each time, it's a geometric sequence. It's a pretty basic concepts that most would've encounter it in high school algebra. So, why do we have to analyze it here?

The reason for this is the definition of a \textbf{limit}, which is the building blocks of derivative:
\begin{equation}
  \odv{f(x)}{x} = \lim_{h \appr 0}\frac{f(x + h) - f(x)}{h},
\end{equation}
i.e., the derivative of $f(x)$ w.r.t. $x$ is just $\frac{f(x + h) - f(x)}{h}$ when $h$ \enquote{approaches} zero. But mathematically, \emph{approaching} is very vague. Some questions may be: How are we approaching the limit? How fast must we approach it? Does this limit always exist? How do we prove that a function approaches something? To answer these and get rid of all ambiguities, one must define \textbf{limit} is using purely mathematical statements.

A first ponder at this defining $\lim_{x \appr a}f(x) = L$ might go something like this. First, let's there be a sequence of number that approaches $a$. If $a = 0$, then that sequence could be
\begin{equation*}
  \ab\{1, \frac{1}{2}, \frac{1}{3}, \frac{1}{4}, \dots\}
\end{equation*}
Then feed in these sequences into $f(x)$ to get another sequence
\begin{equation*}
  \ab\{f(1), f\ab(\frac{1}{2}), f\ab(\frac{1}{3}), f\ab(\frac{1}{4}), \dots\}
\end{equation*}
For $\lim_{x \appr a}f(x) = L$, this sequence must converge to $L$! But now we're stuck. What actually does it mean by convergence? And how do we say that a sequence actually converges to $L$? The problem boils down to defining what a converging sequence means. If we can do that, then it's possible to define what a limit is from the ground up. This chapter shall do just that.

Firstly, let's define what a sequence is
\begin{df}{Sequence}{}
A \textbf{sequence}, denoted $a_n$ is a mapping $a$ from $\mathbb{N}$ to $\mathbb{R}$.
\end{df}
When we want to refer to the whole sequence as a mathematical object, we put a parenthesis around the expression. E.g.,
\begin{gather*}
  (a_n) = \ab(1, \frac{1}{2}, \frac{1}{3}, \dots).
\end{gather*}

\section{Boundedness}

\begin{df}
  A sequence $(a_n)$ is bounded if
  \begin{equation}
    \exists(L \in \mathbb{R})\exists(U \in \mathbb{R})\forall (n \in \mathbb{Z})[L \leq a_n \leq U]
  \end{equation}
\end{df}

\section{Convergence and divergence}
